\begin{proof}[Proof of \cref{obj:lem:paired-tent-full-record-lower}]
\begin{enumerate}
\item \emph{Choice of the mark sequence.}
Fix an admissible tuple \(v=(p,\alpha,\beta,\gamma)\). Define
\[
\vartheta_\triangle=\frac{2q}{2\gamma+q},
\qquad
\kappa_\triangle=\frac1{16},
\qquad
N_{\triangle,n}=2\left\lceil n^{\vartheta_\triangle}\right\rceil,
\qquad
h_{\triangle,n}=N_{\triangle,n}^{-1}
\quad(n\in\mathbb N),
\]
where the reciprocal at zero is defined to be zero. Choose the sequence
\[
L(n)=
\begin{cases}
0,&n=0,\\
N_{\triangle,n}^{\gamma/(p-1)},&n\ge1.
\end{cases}
\]
Admissibility gives \(p-1>0\), \(q>0\), and \(0<\gamma\le1\), hence
\[
\vartheta_\triangle>0,
\qquad
\frac{\gamma}{p-1}>0.
\]
For \(n\ge1\), the ceiling inequality gives
\[
N_{\triangle,n}\ge n^{\vartheta_\triangle}.
\]
Thus \(N_{\triangle,n}\to\infty\), and raising it to the positive power
\(\gamma/(p-1)\) proves \(L(n)\to\infty\).

Now fix \(n\ge2\). Since \(n^{\vartheta_\triangle}\ge1\),
\[
2n^{\vartheta_\triangle}
\le N_{\triangle,n}
<2\bigl(n^{\vartheta_\triangle}+1\bigr)
\le4n^{\vartheta_\triangle}.
\]
In particular, \(N_{\triangle,n}\) is even and at least two. Define
\[
\varepsilon_{\triangle,n}
=h_{\triangle,n}^{\gamma/q}.
\]
Then
\[
0<h_{\triangle,n}<1,
\qquad
0<\varepsilon_{\triangle,n}<1,
\qquad
L(n)=h_{\triangle,n}^{-\gamma/(p-1)}>0.
\]
The identities
\[
\frac{\gamma}{q}-\frac{\gamma}{p-1}=\gamma,
\qquad
\frac{\gamma}{q}-\frac{p\gamma}{p-1}=0
\]
give
\[
\varepsilon_{\triangle,n}L(n)=h_{\triangle,n}^{\gamma},
\qquad
\varepsilon_{\triangle,n}L(n)^p=1.
\]
% lean: CausalSmith.Stat.FinitepHomogeneityDensegamma.tentMagnitude_tendsto_atTop

\item \emph{Paired tents and the finite prior.}
Define the compactly supported tent on the whole real line by
\[
\psi_\triangle(t_\triangle)=
\begin{cases}
2\min\{t_\triangle,1-t_\triangle\},&0\le t_\triangle\le1,\\
0,&\text{otherwise},
\end{cases}
\qquad t_\triangle\in\mathbb R.
\]
Define the number of pairs, their index set, and their sign vectors by
\[
m_{\triangle,n}=\frac{N_{\triangle,n}}2,
\qquad
\mathscr J_{\triangle,n}=\{0,\ldots,m_{\triangle,n}-1\},
\qquad
\varsigma\in\{-1,1\}^{\mathscr J_{\triangle,n}}.
\]
For \(j_\triangle\in\mathscr J_{\triangle,n}\) and \(x\in\mathbb R\), put
\[
\psi_{\triangle,j_\triangle}^{\varsigma}(x)
=
\varsigma_{j_\triangle}
\left[
\psi_\triangle(N_{\triangle,n}x-2j_\triangle)
-\psi_\triangle(N_{\triangle,n}x-(2j_\triangle+1))
\right],
\]
\[
f_{\triangle,\varsigma}(x)
=
\sum_{j_\triangle\in\mathscr J_{\triangle,n}}
\psi_{\triangle,j_\triangle}^{\varsigma}(x).
\]
These functions are continuous. A nonzero paired bump satisfies
\[
\psi_{\triangle,j_\triangle}^{\varsigma}(x)\ne0
\quad\Longrightarrow\quad
2j_\triangle<N_{\triangle,n}x<2j_\triangle+2.
\]
The open intervals in this display are disjoint for distinct pair indices.
Within a pair, the two individual tents also have disjoint nonzero supports.
Consequently,
\[
\left|\psi_{\triangle,j_\triangle}^{\varsigma}(x)\right|\le1,
\qquad
|f_{\triangle,\varsigma}(x)|\le1.
\]

For the Hölder estimate, the identity
\[
\psi_\triangle(t_\triangle)
=\max\{0,2\min(t_\triangle,1-t_\triangle)\}
\]
and the Lipschitz inequalities for minimum and maximum imply
\[
|\psi_\triangle(t_\triangle)-\psi_\triangle(t_\triangle')|
\le2|t_\triangle-t_\triangle'|
\qquad(t_\triangle,t_\triangle'\in\mathbb R).
\]
Applying this to the two tents in a pair gives
\[
\left|
\psi_{\triangle,j_\triangle}^{\varsigma}(x)
-\psi_{\triangle,j_\triangle}^{\varsigma}(x')
\right|
\le4N_{\triangle,n}|x-x'|.
\]
Define the active pairs at a point by
\[
\mathscr A_{\triangle,\varsigma}(x)
=
\left\{
j_\triangle\in\mathscr J_{\triangle,n}:
\psi_{\triangle,j_\triangle}^{\varsigma}(x)\ne0
\right\}.
\]
The support calculation gives
\[
|\mathscr A_{\triangle,\varsigma}(x)|\le1,
\qquad
|\mathscr A_{\triangle,\varsigma}(x)
 \cup\mathscr A_{\triangle,\varsigma}(x')|\le2.
\]
Every summand outside this union vanishes at both points. Therefore
\[
\begin{aligned}
|f_{\triangle,\varsigma}(x)-f_{\triangle,\varsigma}(x')|
&\le
\sum_{j_\triangle\in
\mathscr A_{\triangle,\varsigma}(x)\cup
\mathscr A_{\triangle,\varsigma}(x')}
\left|
\psi_{\triangle,j_\triangle}^{\varsigma}(x)
-\psi_{\triangle,j_\triangle}^{\varsigma}(x')
\right|\\
&\le8N_{\triangle,n}|x-x'|.
\end{aligned}
\]
If \(N_{\triangle,n}|x-x'|\le1\), then
\[
N_{\triangle,n}|x-x'|
\le\bigl(N_{\triangle,n}|x-x'|\bigr)^\gamma.
\]
If \(N_{\triangle,n}|x-x'|>1\), the unit envelope instead gives
\[
|f_{\triangle,\varsigma}(x)-f_{\triangle,\varsigma}(x')|
\le2
\le8\bigl(N_{\triangle,n}|x-x'|\bigr)^\gamma.
\]
Both cases yield
\[
h_{\triangle,n}^{\gamma}
|f_{\triangle,\varsigma}(x)-f_{\triangle,\varsigma}(x')|
\le8|x-x'|^\gamma.
\]

Give every sign vector the same weight, and define
\[
\pi_1
=
2^{-m_{\triangle,n}}
\sum_{\varsigma\in\{-1,1\}^{\mathscr J_{\triangle,n}}}
\delta_{P_\varsigma},
\]
where the laws \(P_\varsigma\) are constructed below. There are exactly
\(2^{m_{\triangle,n}}\) sign vectors, so
\[
2^{-m_{\triangle,n}}>0,
\qquad
\sum_{\varsigma}2^{-m_{\triangle,n}}=1.
\]
The constant sign vector supplies an index, proving that the prior has a
strictly positive weight.
% lean: CausalSmith.Stat.FinitepHomogeneityDensegamma.tentEffect_holderBall

\item \emph{The null law and arm moments.}
Construct \(P_0\) and \(P_\varsigma\) with uniform covariate distribution and
propensity \(1/2\). Their common control kernel and their treated kernels are
\[
Q_{0,P_0}(\,\cdot\mid x)
=
Q_{0,P_\varsigma}(\,\cdot\mid x)=\delta_0,
\]
\[
Q_{1,P_0}(\,\cdot\mid x)
=
(1-\varepsilon_{\triangle,n})\delta_0
+\frac{\varepsilon_{\triangle,n}}2\delta_{L(n)}
+\frac{\varepsilon_{\triangle,n}}2\delta_{-L(n)},
\]
\[
\begin{aligned}
Q_{1,P_\varsigma}(\,\cdot\mid x)
={}&(1-\varepsilon_{\triangle,n})\delta_0\\
&+\frac{\varepsilon_{\triangle,n}}2
 \bigl(1+\kappa_\triangle f_{\triangle,\varsigma}(x)\bigr)\delta_{L(n)}\\
&+\frac{\varepsilon_{\triangle,n}}2
 \bigl(1-\kappa_\triangle f_{\triangle,\varsigma}(x)\bigr)\delta_{-L(n)}.
\end{aligned}
\]
Precisely, these kernels specify the record laws through
\[
P(dx,\{a\},dy)
=\frac12\,\lambda(dx)\,Q_{a,P}(dy\mid x),
\qquad
a\in\{0,1\},
\qquad
P\in\{P_0,P_\varsigma\}.
\]
The displayed masses sum to one. They are nonnegative because
\[
1\pm\kappa_\triangle f_{\triangle,\varsigma}(x)
\ge1-\kappa_\triangle=\frac{15}{16}>0,
\qquad
1-\varepsilon_{\triangle,n}>0.
\]
Continuity of the tents supplies the required kernel measurability.

For every \(x\in[0,1]\), direct integration gives
\[
e_{P_0}(x)=\frac12,
\qquad
m_{0,P_0}(x)=m_{1,P_0}(x)=\tau_{P_0}(x)=0,
\]
and, for \(P=P_0\) or \(P=P_\varsigma\),
\[
Q_{a,P}(\{0,-L(n),L(n)\}\mid x)=1,
\qquad
\int |y|^p\,Q_{a,P}(dy\mid x)
=
\begin{cases}
0,&a=0,\\
\varepsilon_{\triangle,n}L(n)^p=1,&a=1.
\end{cases}
\]
Thus both arm moments are at most one. For \(P_0\), the propensity lies in
\([1/4,3/4]\), its supremum norm is \(1/2\), and its Hölder differences vanish.
The baseline and effect have zero supremum norms and zero Hölder differences.
Together with uniform design and the moment bound \(1\le10\), these facts
verify all the restrictions in \cref{obj:def:model}. Its effect is the constant
zero, so \cref{obj:def:null} gives
\[
P_0\in H_0(v).
\]
% lean: CausalSmith.Stat.FinitepHomogeneityDensegamma.tentLaw_false_inNull

\item \emph{Alternative membership and separation.}
The alternative arm means are
\[
e_{P_\varsigma}(x)=\frac12,
\qquad
m_{0,P_\varsigma}(x)=0,
\]
\[
m_{1,P_\varsigma}(x)=\tau_{P_\varsigma}(x)
=
\varepsilon_{\triangle,n}L(n)\kappa_\triangle
f_{\triangle,\varsigma}(x)
=
\kappa_\triangle h_{\triangle,n}^{\gamma}
f_{\triangle,\varsigma}(x).
\]
The geometric estimates imply
\[
\|\tau_{P_\varsigma}\|_\infty
\le\kappa_\triangle h_{\triangle,n}^{\gamma}
\le\frac1{16}<\frac12,
\]
\[
|\tau_{P_\varsigma}(x)-\tau_{P_\varsigma}(x')|
\le8\kappa_\triangle|x-x'|^\gamma
=\frac12|x-x'|^\gamma
\le20|x-x'|^\gamma.
\]
Thus the effect satisfies its Hölder and supremum restrictions. The constant
propensity and zero baseline satisfy their respective smoothness and bound
restrictions, and the preceding arm calculation supplies the raw moments.
Hence
\[
P_\varsigma\in\mathcal M_v
\qquad\text{for every }\varsigma.
\]

To compute the centered distance, reflection about \(1/2\) gives
\[
\int_0^1\psi_\triangle(t_\triangle)\,dt_\triangle
=
2\int_0^{1/2}2t_\triangle\,dt_\triangle=\frac12,
\]
\[
\int_0^1\psi_\triangle(t_\triangle)^2\,dt_\triangle
=
2\int_0^{1/2}4t_\triangle^2\,dt_\triangle=\frac13.
\]
A change of variables, using the compact support of the tent, yields, for
every integer \(i_\triangle\) with \(0\le i_\triangle<N_{\triangle,n}\),
\[
\int_0^1
\psi_\triangle(N_{\triangle,n}x-i_\triangle)\,\lambda(dx)
=\frac{h_{\triangle,n}}2,
\]
\[
\int_0^1
\psi_\triangle(N_{\triangle,n}x-i_\triangle)^2\,\lambda(dx)
=\frac{h_{\triangle,n}}3.
\]
The opposite signs cancel the two equal areas in each pair. Disjoint nonzero
supports eliminate the cross terms in the square. Consequently,
\[
\int f_{\triangle,\varsigma}\,d\lambda
=
\sum_{j_\triangle}\varsigma_{j_\triangle}
\left(\frac{h_{\triangle,n}}2-\frac{h_{\triangle,n}}2\right)=0,
\]
\[
\int f_{\triangle,\varsigma}^2\,d\lambda
=
m_{\triangle,n}\frac{2h_{\triangle,n}}3=\frac13.
\]
It follows that
\[
\int\tau_{P_\varsigma}\,d\lambda=0,
\qquad
d(P_\varsigma)
=\frac{\kappa_\triangle h_{\triangle,n}^{\gamma}}{\sqrt3}.
\]
The upper rank bound supplies
\[
h_{\triangle,n}^{\gamma}
\ge(4n^{\vartheta_\triangle})^{-\gamma}
=4^{-\gamma}n^{-E_0},
\]
because \(\gamma\vartheta_\triangle=E_0\). Therefore
\[
d(P_\varsigma)\ge
\frac{4^{-\gamma}}{16\sqrt3}n^{-E_0}
=c_v^{\mathrm e}n^{-E_0}.
\]

For the strict upper comparison, expand the centered square and use the
effect envelope:
\[
\begin{aligned}
d(P_\varsigma)^2
&=
\int\tau_{P_\varsigma}^2\,d\lambda
-\left(\int\tau_{P_\varsigma}\,d\lambda\right)^2\\
&\le\int\tau_{P_\varsigma}^2\,d\lambda
\le\frac1{16^2}.
\end{aligned}
\]
Since \(\sqrt2<2\),
\[
d(P_\varsigma)\le\frac1{16}
<\frac1{8\sqrt2}=d_0.
\]
% lean: CausalSmith.Stat.FinitepHomogeneityDensegamma.paired_effect_distance_lower

\item \emph{Comparison of the complete-record laws.}
Define the one-record likelihood on
\([0,1]\times\{0,1\}\times\mathbb R\) by
\[
\ell_{\triangle,\varsigma}(x,a,y)=
\begin{cases}
1+\kappa_\triangle f_{\triangle,\varsigma}(x),
&a=1,\ y=L(n),\\
1-\kappa_\triangle f_{\triangle,\varsigma}(x),
&a=1,\ y=-L(n),\\
1,&\text{otherwise}.
\end{cases}
\]
Comparison of the conditional atom masses proves
\[
P_\varsigma\ll P_0,
\qquad
\frac{dP_\varsigma}{dP_0}=\ell_{\triangle,\varsigma}
\quad P_0\text{-almost surely},
\qquad
\frac{15}{16}\le\ell_{\triangle,\varsigma}\le\frac{17}{16}.
\]
At each covariate value, summing over the control record and the three
treated atoms gives
\[
\begin{aligned}
&\frac12+\frac{1-\varepsilon_{\triangle,n}}2\\
&\quad+\frac{\varepsilon_{\triangle,n}}4
\left[
(1+\kappa_\triangle f_{\triangle,\varsigma}(x))
(1+\kappa_\triangle f_{\triangle,\varsigma'}(x))
+
(1-\kappa_\triangle f_{\triangle,\varsigma}(x))
(1-\kappa_\triangle f_{\triangle,\varsigma'}(x))
\right]\\
&\hspace{2cm}
=1+\frac{\varepsilon_{\triangle,n}\kappa_\triangle^2}{2}
f_{\triangle,\varsigma}(x)f_{\triangle,\varsigma'}(x).
\end{aligned}
\]
The same support calculation used for the squared tents gives
\[
\begin{aligned}
\int f_{\triangle,\varsigma}f_{\triangle,\varsigma'}\,d\lambda
&=
\sum_{j_\triangle}
\varsigma_{j_\triangle}\varsigma'_{j_\triangle}
\left(\frac{h_{\triangle,n}}3+\frac{h_{\triangle,n}}3\right)\\
&=\frac{2h_{\triangle,n}}3
\sum_{j_\triangle}\varsigma_{j_\triangle}\varsigma'_{j_\triangle}.
\end{aligned}
\]
Define the overlap coefficient by
\[
\varrho_{\triangle,n}
=\frac{\varepsilon_{\triangle,n}\kappa_\triangle^2}{6}.
\]
Integration over the entire record now yields
\[
\int\ell_{\triangle,\varsigma}\ell_{\triangle,\varsigma'}\,dP_0
=
1+\frac{\varrho_{\triangle,n}}{m_{\triangle,n}}
\sum_{j_\triangle}\varsigma_{j_\triangle}\varsigma'_{j_\triangle}.
\]
Here \(m_{\triangle,n}\ge1\), \(0\le\varrho_{\triangle,n}\le1\), and
\[
\left|
\sum_{j_\triangle}\varsigma_{j_\triangle}\varsigma'_{j_\triangle}
\right|\le m_{\triangle,n}.
\]
Thus every overlap base is nonnegative.

For probability laws \(\mu_{\rm alt}\ll\mu_{\rm ref}\) with square-integrable
density, define the directed chi-square divergence by
\[
\chi^2(\mu_{\rm alt}\Vert\mu_{\rm ref})
=
\int
\left(\frac{d\mu_{\rm alt}}{d\mu_{\rm ref}}-1\right)^2
\,d\mu_{\rm ref}.
\]
The complete-record mixture and its density are
\[
\overline P_1^{(n)}
=
2^{-m_{\triangle,n}}\sum_\varsigma P_\varsigma^{\otimes n},
\]
\[
G_\triangle(\mathcal D_n)
=
\frac{d\overline P_1^{(n)}}{dP_0^{\otimes n}}(\mathcal D_n)
=
2^{-m_{\triangle,n}}\sum_\varsigma
\prod_{i_{\rm rec}=1}^n
\ell_{\triangle,\varsigma}(X_{i_{\rm rec}},A_{i_{\rm rec}},Y_{i_{\rm rec}}).
\]
These are bounded finite densities, so expansion of their square and
independence of the records give
\[
\begin{aligned}
1+\chi^2(\overline P_1^{(n)}\Vert P_0^{\otimes n})
&=\int G_\triangle^2\,dP_0^{\otimes n}\\
&=
2^{-2m_{\triangle,n}}
\sum_\varsigma\sum_{\varsigma'}
\left(
1+\frac{\varrho_{\triangle,n}}{m_{\triangle,n}}
\sum_{j_\triangle}\varsigma_{j_\triangle}\varsigma'_{j_\triangle}
\right)^n.
\end{aligned}
\]
Nonnegativity of the bases and \(1+t_\triangle\le e^{t_\triangle}\) imply
\[
\left(
1+\frac{\varrho_{\triangle,n}}{m_{\triangle,n}}
\sum_{j_\triangle}\varsigma_{j_\triangle}\varsigma'_{j_\triangle}
\right)^n
\le
\exp\left(
\frac{n\varrho_{\triangle,n}}{m_{\triangle,n}}
\sum_{j_\triangle}\varsigma_{j_\triangle}\varsigma'_{j_\triangle}
\right).
\]
For each fixed \(\varsigma\), the sum over \(\varsigma'\) factors coordinate
by coordinate:
\[
\begin{aligned}
&2^{-m_{\triangle,n}}\sum_{\varsigma'}
\exp\left(
\frac{n\varrho_{\triangle,n}}{m_{\triangle,n}}
\sum_{j_\triangle}\varsigma_{j_\triangle}\varsigma'_{j_\triangle}
\right)\\
&\qquad=
\prod_{j_\triangle}
\cosh\left(
\frac{n\varrho_{\triangle,n}}{m_{\triangle,n}}\varsigma_{j_\triangle}
\right)
\le
\exp\left(\frac{(n\varrho_{\triangle,n})^2}{2m_{\triangle,n}}\right).
\end{aligned}
\]
The last inequality follows from
\[
\cosh(t_\triangle)
=
\sum_{i_{\rm exp}=0}^{\infty}
\frac{t_\triangle^{2i_{\rm exp}}}{(2i_{\rm exp})!}
\le
\sum_{i_{\rm exp}=0}^{\infty}
\frac{(t_\triangle^2/2)^{i_{\rm exp}}}{i_{\rm exp}!}
=e^{t_\triangle^2/2},
\]
using \((2i_{\rm exp})!\ge2^{i_{\rm exp}}i_{\rm exp}!\).
Averaging also over \(\varsigma\) therefore gives
\[
1+\chi^2(\overline P_1^{(n)}\Vert P_0^{\otimes n})
\le
\exp\left(
\frac{\kappa_\triangle^4}{36}
n^2\varepsilon_{\triangle,n}^2h_{\triangle,n}
\right).
\]
Indeed, \(2m_{\triangle,n}=N_{\triangle,n}=h_{\triangle,n}^{-1}\).
Moreover,
\[
h_{\triangle,n}\le n^{-\vartheta_\triangle},
\qquad
\vartheta_\triangle\left(\frac{2\gamma}{q}+1\right)=2,
\]
so
\[
n^2\varepsilon_{\triangle,n}^2h_{\triangle,n}
=
n^2h_{\triangle,n}^{2\gamma/q+1}
\le n^2n^{-2}=1.
\]
Consequently,
\[
1+\chi^2(\overline P_1^{(n)}\Vert P_0^{\otimes n})
\le e^{\kappa_\triangle^4/36}
\le\frac1{1-\kappa_\triangle^4/36}<2.
\]
The middle inequality follows by comparing the exponential and geometric
series; the last uses \(\kappa_\triangle=1/16\).

Since \(\int(G_\triangle-1)\,dP_0^{\otimes n}=0\), its positive and negative
parts have equal integrals. The event \(\{G_\triangle\ge1\}\) then gives
\[
\operatorname{TV}(P_0^{\otimes n},\overline P_1^{(n)})
=
\frac12\int|G_\triangle-1|\,dP_0^{\otimes n}.
\]
Cauchy--Schwarz under the probability measure \(P_0^{\otimes n}\) yields
\[
\operatorname{TV}(P_0^{\otimes n},\overline P_1^{(n)})
\le
\frac12
\sqrt{\chi^2(\overline P_1^{(n)}\Vert P_0^{\otimes n})}
<\frac12.
\]
% lean: CausalSmith.Stat.FinitepHomogeneityDensegamma.tentMixture_tv_lt_half

\item \emph{Supplied-propensity and original-record risks.}
By \cref{obj:lem:causal-nonempty},
\[
d_0\le D_v.
\]
Define the separation used here by
\[
r_\triangle=c_v^{\mathrm e}n^{-E_0}.
\]
The positive prior weights and the established placement give
\[
0<r_\triangle\le d(P_\varsigma)<d_0\le D_v.
\]

Take any supplied-propensity test satisfying the size constraint in
\cref{obj:def:oracle-risk}. Define its section at the common propensity and
its seed average by
\[
\varpi_{\mathrm{seed}}\in[0,1],
\qquad
\phi_{1/2}(\mathcal D_n,\varpi_{\mathrm{seed}})
=
\phi^{\mathrm e}(\mathcal D_n,\varpi_{\mathrm{seed}};\,e\equiv1/2),
\qquad
\overline\phi_{1/2}(\mathcal D_n)
=
\int_0^1\phi_{1/2}(\mathcal D_n,\varpi_{\mathrm{seed}})\,\lambda(d\varpi_{\mathrm{seed}}).
\]
This is a measurable function with values in \([0,1]\). Integrating its level
sets shows that its expectation difference under two probability laws is
bounded by their total variation. Thus
\[
\left|
2^{-m_{\triangle,n}}\sum_\varsigma
\mathsf R_n^{\mathrm e}(P_\varsigma,\phi^{\mathrm e})
-
\mathsf R_n^{\mathrm e}(P_0,\phi^{\mathrm e})
\right|
\le
\operatorname{TV}(P_0^{\otimes n},\overline P_1^{(n)})
<\frac12.
\]
Since \(P_0\in H_0(v)\),
\[
2^{-m_{\triangle,n}}\sum_\varsigma
\mathsf R_n^{\mathrm e}(P_\varsigma,\phi^{\mathrm e})
<\frac1{10}+\frac12=\frac35.
\]
If every positive-weight alternative had type-II error strictly below
\(2/5\), normalization of the weights would instead imply
\[
2^{-m_{\triangle,n}}\sum_\varsigma
\mathsf R_n^{\mathrm e}(P_\varsigma,\phi^{\mathrm e})
\ge
2^{-m_{\triangle,n}}\sum_\varsigma\frac35
=\frac35.
\]
Hence some support law has type-II error at least \(2/5\). Its model
membership and separation imply
\[
\sup_{\substack{P\in\mathcal M_v\\d(P)\ge r_\triangle}}
\bigl\{1-\mathsf R_n^{\mathrm e}(P,\phi^{\mathrm e})\bigr\}
\ge\frac25.
\]
The constant-zero test makes the admissible test class nonempty, and all
rejection probabilities lie in \([0,1]\). Taking the infimum therefore gives
\[
B_n^{\mathrm e}(v,r_\triangle)\ge\frac25.
\]

For an original-record test \(\phi\), use instead
\[
\overline\phi(\mathcal D_n)
=\int_0^1\phi(\mathcal D_n,\varpi_{\mathrm{seed}})\,\lambda(d\varpi_{\mathrm{seed}}).
\]
The identical expectation comparison gives
\[
2^{-m_{\triangle,n}}\sum_\varsigma
\mathsf R_n(P_\varsigma,\phi)<\frac35.
\]
The same weighted-sum contradiction supplies a support law with type-II
error at least \(2/5\). Taking the supremum and then the infimum in
\cref{obj:def:testing-risk} proves
\[
B_n(v,r_\triangle)\ge\frac25.
\]
This completes the first witness package.
% lean: CausalSmith.Stat.FinitepHomogeneityDensegamma.tent_canonical_lower_receipt

\item \emph{The signed-binary construction.}
Fix an admissible \(w=(\alpha,\beta,\gamma)\) and \(n\ge2\). Define
\[
\vartheta_\triangle^{\mathrm b}=\frac2{4\gamma+1},
\qquad
N_{\triangle,n}^{\mathrm b}
=2\left\lceil n^{\vartheta_\triangle^{\mathrm b}}\right\rceil,
\qquad
h_{\triangle,n}^{\mathrm b}=(N_{\triangle,n}^{\mathrm b})^{-1},
\]
\[
m_{\triangle,n}^{\mathrm b}=\frac{N_{\triangle,n}^{\mathrm b}}2,
\qquad
\mathscr J_{\triangle,n}^{\mathrm b}
=\{0,\ldots,m_{\triangle,n}^{\mathrm b}-1\},
\qquad
\delta_{\triangle,n}^{\mathrm b}
=\kappa_\triangle(h_{\triangle,n}^{\mathrm b})^\gamma.
\]
For \(\varsigma\in\{-1,1\}^{\mathscr J_{\triangle,n}^{\mathrm b}}\), define
\[
f_{\triangle,\varsigma}^{\mathrm b}(x)
=
\sum_{j_\triangle\in\mathscr J_{\triangle,n}^{\mathrm b}}
\varsigma_{j_\triangle}
\left[
\psi_\triangle(N_{\triangle,n}^{\mathrm b}x-2j_\triangle)
-\psi_\triangle(N_{\triangle,n}^{\mathrm b}x-(2j_\triangle+1))
\right],
\qquad x\in\mathbb R.
\]
These are the preceding paired-tent constructions evaluated at moment order
two, for which \(q=1/2\). In particular,
\[
2n^{\vartheta_\triangle^{\mathrm b}}
\le N_{\triangle,n}^{\mathrm b}
\le4n^{\vartheta_\triangle^{\mathrm b}},
\qquad
0<h_{\triangle,n}^{\mathrm b}<1,
\]
\[
|f_{\triangle,\varsigma}^{\mathrm b}(x)|\le1,
\qquad
(h_{\triangle,n}^{\mathrm b})^\gamma
|f_{\triangle,\varsigma}^{\mathrm b}(x)
-f_{\triangle,\varsigma}^{\mathrm b}(x')|
\le8|x-x'|^\gamma.
\]

To establish legality through the signed-binary conversion, first construct
deterministic record laws with uniform design, propensity \(1/2\), baseline
zero, and outcomes
\[
Y=0
\quad\text{under }P_{\mathrm{det},0}^{\mathrm b},
\qquad
Y=A\delta_{\triangle,n}^{\mathrm b}
f_{\triangle,\varsigma}^{\mathrm b}(X)
\quad\text{under }P_{\mathrm{det},\varsigma}^{\mathrm b}.
\]
Their effects are respectively zero and
\(\delta_{\triangle,n}^{\mathrm b}f_{\triangle,\varsigma}^{\mathrm b}\).
The displayed geometric estimates give
\[
\left\|
\delta_{\triangle,n}^{\mathrm b}f_{\triangle,\varsigma}^{\mathrm b}
\right\|_\infty\le\frac1{16},
\]
\[
\left|
\delta_{\triangle,n}^{\mathrm b}f_{\triangle,\varsigma}^{\mathrm b}(x)
-\delta_{\triangle,n}^{\mathrm b}f_{\triangle,\varsigma}^{\mathrm b}(x')
\right|
\le\frac12|x-x'|^\gamma.
\]
The control second moments are zero, and the treated second moments satisfy
\[
\bigl|\delta_{\triangle,n}^{\mathrm b}
f_{\triangle,\varsigma}^{\mathrm b}(x)\bigr|^2
\le\frac1{16^2}\le10.
\]
Together with the constant propensity and zero baseline, these inequalities
put both deterministic laws in
\(\mathcal M_{(2,\alpha,\beta,\gamma)}\); the first has constant zero effect.

Apply the conversion in \cref{obj:prop:bounded-submodel-comparison}. It
preserves the design, propensity, means, effect, and null constancy. The
resulting laws \(P_0^{\mathrm b}\) and \(P_\varsigma^{\mathrm b}\) have kernels
\[
Q_{a,P_0^{\mathrm b}}(\,\cdot\mid x)
=\frac12\delta_1+\frac12\delta_{-1}
\qquad(a\in\{0,1\}),
\]
\[
Q_{0,P_\varsigma^{\mathrm b}}(\,\cdot\mid x)
=\frac12\delta_1+\frac12\delta_{-1},
\]
\[
Q_{1,P_\varsigma^{\mathrm b}}(\,\cdot\mid x)
=
\frac{1+\delta_{\triangle,n}^{\mathrm b}
f_{\triangle,\varsigma}^{\mathrm b}(x)}2\delta_1
+
\frac{1-\delta_{\triangle,n}^{\mathrm b}
f_{\triangle,\varsigma}^{\mathrm b}(x)}2\delta_{-1}.
\]
Their propensity is identically \(1/2\), and
\[
P_0^{\mathrm b}\in H_0^{\mathrm{bin}}(w),
\qquad
P_\varsigma^{\mathrm b}\in\mathcal M_w^{\mathrm{bin}},
\qquad
\tau_{P_\varsigma^{\mathrm b}}
=\delta_{\triangle,n}^{\mathrm b}f_{\triangle,\varsigma}^{\mathrm b}.
\]
Define the binary prior by
\[
\pi_1^{\mathrm b}
=
2^{-m_{\triangle,n}^{\mathrm b}}
\sum_{\varsigma\in\{-1,1\}^{\mathscr J_{\triangle,n}^{\mathrm b}}}
\delta_{P_\varsigma^{\mathrm b}}.
\]
Again,
\[
2^{-m_{\triangle,n}^{\mathrm b}}>0,
\qquad
\sum_\varsigma2^{-m_{\triangle,n}^{\mathrm b}}=1,
\]
and the constant sign vector gives a positive-weight index.
% lean: CausalSmith.Stat.FinitepHomogeneityDensegamma.binaryTentLaw_inBinaryModel

\item \emph{Signed-binary separation.}
The equal-area cancellation and the disjoint-support square calculation give
\[
\int f_{\triangle,\varsigma}^{\mathrm b}\,d\lambda=0,
\qquad
\int(f_{\triangle,\varsigma}^{\mathrm b})^2\,d\lambda=\frac13.
\]
Consequently,
\[
d(P_\varsigma^{\mathrm b})
=\frac{\kappa_\triangle(h_{\triangle,n}^{\mathrm b})^\gamma}{\sqrt3}.
\]
Using the upper rank bound and
\(\gamma\vartheta_\triangle^{\mathrm b}=2\gamma/(4\gamma+1)\),
\[
\begin{aligned}
d(P_\varsigma^{\mathrm b})
&\ge
\frac{\kappa_\triangle}{\sqrt3}
(4n^{\vartheta_\triangle^{\mathrm b}})^{-\gamma}\\
&=
c_w^{\mathrm b}n^{-2\gamma/(4\gamma+1)}.
\end{aligned}
\]
The effect envelope also gives
\[
\begin{aligned}
d(P_\varsigma^{\mathrm b})^2
&=
\int\tau_{P_\varsigma^{\mathrm b}}^2\,d\lambda
-\left(\int\tau_{P_\varsigma^{\mathrm b}}\,d\lambda\right)^2\\
&\le\frac1{16^2},
\end{aligned}
\qquad
d(P_\varsigma^{\mathrm b})\le\frac1{16}<d_0.
\]
% lean: CausalSmith.Stat.FinitepHomogeneityDensegamma.binaryTentLaw_true_distance_lt_d0

\item \emph{Signed-binary complete-record comparison.}
Define, on the whole record space,
\[
\ell_{\triangle,\varsigma}^{\mathrm b}(x,a,y)=
\begin{cases}
1+\delta_{\triangle,n}^{\mathrm b}
f_{\triangle,\varsigma}^{\mathrm b}(x),&a=1,\ y=1,\\
1-\delta_{\triangle,n}^{\mathrm b}
f_{\triangle,\varsigma}^{\mathrm b}(x),&a=1,\ y=-1,\\
1,&\text{otherwise}.
\end{cases}
\]
The kernel probabilities imply
\[
P_\varsigma^{\mathrm b}\ll P_0^{\mathrm b},
\qquad
\frac{dP_\varsigma^{\mathrm b}}{dP_0^{\mathrm b}}
=\ell_{\triangle,\varsigma}^{\mathrm b}
\quad P_0^{\mathrm b}\text{-almost surely}.
\]
Under the null, each treatment--sign combination has conditional probability
\(1/4\). Summing the two control terms and the two treated terms yields
\[
\begin{aligned}
&\frac12+\frac14
\left[
(1+\delta_{\triangle,n}^{\mathrm b}f_{\triangle,\varsigma}^{\mathrm b}(x))
(1+\delta_{\triangle,n}^{\mathrm b}f_{\triangle,\varsigma'}^{\mathrm b}(x))
\right.\\[-2pt]
&\hspace{3cm}\left.
+
(1-\delta_{\triangle,n}^{\mathrm b}f_{\triangle,\varsigma}^{\mathrm b}(x))
(1-\delta_{\triangle,n}^{\mathrm b}f_{\triangle,\varsigma'}^{\mathrm b}(x))
\right]\\
&\qquad=
1+\frac{(\delta_{\triangle,n}^{\mathrm b})^2}{2}
f_{\triangle,\varsigma}^{\mathrm b}(x)
f_{\triangle,\varsigma'}^{\mathrm b}(x).
\end{aligned}
\]
Define
\[
\varepsilon_{\triangle,n}^{\mathrm b}
=(h_{\triangle,n}^{\mathrm b})^{2\gamma},
\qquad
\varrho_{\triangle,n}^{\mathrm b}
=\frac{\varepsilon_{\triangle,n}^{\mathrm b}\kappa_\triangle^2}{6}.
\]
Then
\[
(\delta_{\triangle,n}^{\mathrm b})^2
=\kappa_\triangle^2\varepsilon_{\triangle,n}^{\mathrm b},
\]
and integration of the paired squared tents gives
\[
\int f_{\triangle,\varsigma}^{\mathrm b}
f_{\triangle,\varsigma'}^{\mathrm b}\,d\lambda
=
\frac{2h_{\triangle,n}^{\mathrm b}}3
\sum_{j_\triangle}
\varsigma_{j_\triangle}\varsigma'_{j_\triangle}.
\]
Thus the complete one-record overlap is
\[
\int
\ell_{\triangle,\varsigma}^{\mathrm b}
\ell_{\triangle,\varsigma'}^{\mathrm b}\,dP_0^{\mathrm b}
=
1+\frac{\varrho_{\triangle,n}^{\mathrm b}}{m_{\triangle,n}^{\mathrm b}}
\sum_{j_\triangle}\varsigma_{j_\triangle}\varsigma'_{j_\triangle}.
\]
These bases are nonnegative, since
\[
0\le\varrho_{\triangle,n}^{\mathrm b}\le1,
\qquad
\left|\sum_{j_\triangle}
\varsigma_{j_\triangle}\varsigma'_{j_\triangle}\right|
\le m_{\triangle,n}^{\mathrm b}.
\]

Define the mixture and its density by
\[
\overline P_{1,\mathrm b}^{(n)}
=
2^{-m_{\triangle,n}^{\mathrm b}}
\sum_\varsigma(P_\varsigma^{\mathrm b})^{\otimes n},
\]
\[
G_\triangle^{\mathrm b}(\mathcal D_n)
=
\frac{d\overline P_{1,\mathrm b}^{(n)}}
 {d(P_0^{\mathrm b})^{\otimes n}}(\mathcal D_n)
=
2^{-m_{\triangle,n}^{\mathrm b}}
\sum_\varsigma\prod_{i_{\rm rec}=1}^n
\ell_{\triangle,\varsigma}^{\mathrm b}
(X_{i_{\rm rec}},A_{i_{\rm rec}},Y_{i_{\rm rec}}).
\]
Expansion and independence give
\[
\begin{aligned}
1+\chi^2\!\left(
\overline P_{1,\mathrm b}^{(n)}
\Vert(P_0^{\mathrm b})^{\otimes n}\right)
={}&2^{-2m_{\triangle,n}^{\mathrm b}}
\sum_\varsigma\sum_{\varsigma'}
\left(
1+\frac{\varrho_{\triangle,n}^{\mathrm b}}{m_{\triangle,n}^{\mathrm b}}
\sum_{j_\triangle}
\varsigma_{j_\triangle}\varsigma'_{j_\triangle}
\right)^n.
\end{aligned}
\]
For fixed \(\varsigma\), the exponential comparison and sign factorization
already established give
\[
\begin{aligned}
&2^{-m_{\triangle,n}^{\mathrm b}}\sum_{\varsigma'}
\left(
1+\frac{\varrho_{\triangle,n}^{\mathrm b}}{m_{\triangle,n}^{\mathrm b}}
\sum_{j_\triangle}
\varsigma_{j_\triangle}\varsigma'_{j_\triangle}
\right)^n\\
&\quad\le
\prod_{j_\triangle}
\cosh\left(
\frac{n\varrho_{\triangle,n}^{\mathrm b}}
 {m_{\triangle,n}^{\mathrm b}}\varsigma_{j_\triangle}
\right)
\le
\exp\left(
\frac{(n\varrho_{\triangle,n}^{\mathrm b})^2}
 {2m_{\triangle,n}^{\mathrm b}}\right).
\end{aligned}
\]
Averaging over the remaining signs and simplifying the exponent yields
\[
1+\chi^2\!\left(
\overline P_{1,\mathrm b}^{(n)}
\Vert(P_0^{\mathrm b})^{\otimes n}\right)
\le
\exp\left(
\frac{\kappa_\triangle^4}{36}
n^2(h_{\triangle,n}^{\mathrm b})^{4\gamma+1}
\right).
\]
Since
\[
h_{\triangle,n}^{\mathrm b}
\le n^{-\vartheta_\triangle^{\mathrm b}},
\qquad
\vartheta_\triangle^{\mathrm b}(4\gamma+1)=2,
\]
we have
\[
n^2(h_{\triangle,n}^{\mathrm b})^{4\gamma+1}\le1,
\]
and therefore
\[
1+\chi^2\!\left(
\overline P_{1,\mathrm b}^{(n)}
\Vert(P_0^{\mathrm b})^{\otimes n}\right)
\le e^{\kappa_\triangle^4/36}<2.
\]
Applying the zero-integral density-deviation identity and
Cauchy--Schwarz to \(G_\triangle^{\mathrm b}\) gives
\[
\begin{aligned}
\operatorname{TV}\!\left(
(P_0^{\mathrm b})^{\otimes n},\overline P_{1,\mathrm b}^{(n)}
\right)
&=
\frac12\int|G_\triangle^{\mathrm b}-1|\,d(P_0^{\mathrm b})^{\otimes n}\\
&\le
\frac12\sqrt{\chi^2\!\left(
\overline P_{1,\mathrm b}^{(n)}
\Vert(P_0^{\mathrm b})^{\otimes n}\right)}
<\frac12.
\end{aligned}
\]
% lean: CausalSmith.Stat.FinitepHomogeneityDensegamma.binaryTentMixture_tv_lt_half

\item \emph{Bounded-model placement and risk.}
The inclusion supplied by \cref{obj:prop:bounded-submodel-comparison} gives
\[
P_\varsigma^{\mathrm b}\in\mathcal M_w^{\mathrm b},
\qquad
P_0^{\mathrm b}\in H_0^{\mathrm b}(w).
\]
For the null, this uses its bounded-model membership together with the
preserved constant zero effect. The same comparison, together with
\cref{obj:lem:causal-nonempty} at moment order two, supplies
\[
D_w^{\mathrm b}=D_{(2,\alpha,\beta,\gamma)}\ge d_0.
\]
Define
\[
r_\triangle^{\mathrm b}
=c_w^{\mathrm b}n^{-2\gamma/(4\gamma+1)}.
\]
The binary alternatives satisfy
\[
0<r_\triangle^{\mathrm b}
\le d(P_\varsigma^{\mathrm b})<d_0\le D_w^{\mathrm b}.
\]

For any test satisfying the bounded-null size constraint in
\cref{obj:def:bounded-risk}, average its public seed as above. The
complete-record total-variation bound gives
\[
\left|
2^{-m_{\triangle,n}^{\mathrm b}}\sum_\varsigma
\mathsf R_n(P_\varsigma^{\mathrm b},\phi)
-\mathsf R_n(P_0^{\mathrm b},\phi)
\right|<\frac12,
\]
hence
\[
2^{-m_{\triangle,n}^{\mathrm b}}\sum_\varsigma
\mathsf R_n(P_\varsigma^{\mathrm b},\phi)<\frac35.
\]
If all positive-weight support laws had type-II error below \(2/5\), their
weighted rejection sum would be at least
\[
2^{-m_{\triangle,n}^{\mathrm b}}\sum_\varsigma\frac35=\frac35,
\]
a contradiction. Thus a support law has type-II error at least \(2/5\).
All support laws belong to the bounded model at separation
\(r_\triangle^{\mathrm b}\), so
\[
\sup_{\substack{P\in\mathcal M_w^{\mathrm b}\\
d(P)\ge r_\triangle^{\mathrm b}}}
\{1-\mathsf R_n(P,\phi)\}\ge\frac25.
\]
The constant-zero test is admissible, and taking the infimum over the
bounded-null level class proves
\[
B_n^{\mathrm b}(w,r_\triangle^{\mathrm b})\ge\frac25.
\]
The constructed signed-binary null, normalized positive-weight prior,
placement, and complete-record comparison establish the second witness
package.
% lean: CausalSmith.Stat.FinitepHomogeneityDensegamma.binary_tent_canonical_lower_receipt
\end{enumerate}
\end{proof}
